% GPT-6 Astra Pro; effort not exposed; /root/replacement; 22 November 2026.
\section[HD-F1-R35: unrestricted endpoint ranks at small forward Hom]{HD-F1-R35: three exceptional character coefficients for every adjacent object with at most three invariant forward maps}
\label{sec:followup-rm-unrestricted-low-h-character}

Retain the genuine RM fourfold, rational algebraic $B$,
balanced vector $v=v_1v_2$, $v_i=1-qA_i$, $q>0$,
and a specified coherent faithful finite point action.
Every actual gluable weakly semiregular lifted object with
\[
 \begin{gathered}
 H^0(E)=A,\quad H^1(E)=T,\quad H^j(E)=0\ (j\ne0,1),\\
 \operatorname{ch}(E)=Ne^{-B}v,\quad N\ne0,\qquad
 h=\dim\operatorname{Hom}_G(A,T)\le3,
 \end{gathered}
 \tag{F1R35.1}
\]
where $A,T$ are positive-rank coherent sheaves, satisfies
the following restrictions. There is no torsion-free,
semistability, endpoint-rank or endomorphism-dimension
assumption. All original coupling and coherent-defect
maps are retained.

One has $h=2$ or $3$. Under the actual cohomological
shear $z_A=e^B\operatorname{ch}(A)$,
$z_T=e^B\operatorname{ch}(T)$, the whole characters have
the common product form
\[
 \begin{gathered}
 z_T=a v+c p_1p_2,\qquad z_A=(a+N)v+c p_1p_2,\\
 c\ne0,\quad p_i\notin\mathbf C v_i,\qquad
 a\in\{0,-N,-N/2\}.
 \end{gathered}
 \tag{F1R35.2}
\]
The factors and rationality are precisely the $F$-defined,
factor-conjugate Hodge tensors of R26, with $a,c\in\mathbf Q$.
The shear changes no object: its actual source and target
isomorphisms are recalled below.

The remaining possibilities have the following sharper
operator restrictions. Set
$\alpha_i=\Phi_i^{-1}(p_i)$,
$\xi_i=\Phi_i^{-1}(v_i)$ and
$P_i=L_{\xi_i}^{-1}L_{\alpha_i}$ as in R34.
\[
\begin{array}{c|c|c|c}
 a&h&\text{roots of each }P_i&\text{required equality}\\ \hline
 0&2\text{ or }3&0,\lambda_i\ (\lambda_i\ne0)
          &c\lambda_1\lambda_2=-N\rule{0pt}{2.6ex}\\
 -N&2\text{ or }3&0,\lambda_i\ (\lambda_i\ne0)
          &c\lambda_1\lambda_2=N\\
 -N/2&3&\lambda_i,-\lambda_i\ (\lambda_i\ne0)
          &c^2\lambda_1^2\lambda_2^2=N^2/4.
\end{array}
\tag{F1R35.3}
\]
Every operator in this table is semisimple, with each
listed eigenspace of dimension two. In the first two
rows the endpoint with zero balanced coefficient has
entire character
\[
          \operatorname{ch}(F)=\operatorname{rk}(F)e^{\ell-B},
               \qquad \ell\in H^{1,1}(X)\cap H^2(X,\mathbf Q).
 \tag{F1R35.4}
\]
This is a necessary identity for the existing sheaf;
no bundle, sheaf, extension or semiregular object is
asserted to realize any row. At $N=1$, the three
coefficients in (F1R35.2) are exactly $0,-1,-1/2$.
R34 already excludes the rank-$(3,2)$ equal-slope
semistable TF gate. The present restrictions concern
the genuinely wider endpoint class (F1R35.1).

\paragraph{1. The general triangle comparison retains every kernel.}
Use the original lifted triangle
\[
 A\xrightarrow{i}E\xrightarrow{p}C=T[-1]
             \xrightarrow{d=\delta[-1]}A[1],
                    \qquad \delta:T\longrightarrow A[2].
 \tag{F1R35.5}
\]
Put $W=\operatorname{Hom}_G(A,T)$ and define
\[
 \begin{gathered}
 \nu(f)=\delta f,\quad\mu(f)=f[2]\delta,\quad
 J(f)=\delta[2]f[2]\delta,\\
 D:\operatorname{Ext}^2_G(A,A)\oplus\operatorname{Ext}^2_G(T,T)
       \longrightarrow\operatorname{Ext}^4_G(T,A),\\
 D(a',t')=a'[2]\delta-\delta[2]t'.
 \end{gathered}
 \tag{F1R35.6}
\]
Let $r_\nu,r_\mu,r_J,r_D$ be their ranks.
No injectivity is assumed. Invariant CY4 Serre duality
identifies the target of $D$ with $W^*$, so $r_D\le h$.

Let $\mathcal M$ be the vector space of actual triples
$(a',e,t')$ giving a morphism from (F1R35.5) to its
shift by two:
\[
 ei=i[2]a',\qquad p[2]e=t'[-1]p,\qquad
                        a'[2]\delta=\delta[2]t'.
 \tag{F1R35.7}
\]
Its endpoint projection is onto $\ker D$. Indeed,
twice rotate the original triangle to
\[
 C\xrightarrow{d}A[1]\xrightarrow{-i[1]}E[1]
                                      \xrightarrow{-p[1]}C[1].
 \tag{F1R35.8}
\]
The prescribed first two vertical arrows $t'[-1]$ and
$a'[1]$ commute exactly when $D(a',t')=0$.
TR3 gives the third arrow; shifting back gives (F1R35.7).
The two minus signs occur in both rows and cancel.
An ordinary filler can be averaged over the specified
finite action, retaining its invariant endpoints.

For the middle projection $\mathcal M\to\operatorname{Ext}^2_G(E,E)$,
the exact Hom predecessors are
\[
 \begin{split}
 \operatorname{Hom}_G(A,C[1])=W
       &\longrightarrow\operatorname{Ext}^2_G(A,A)
         \xrightarrow{i[2](-)}\operatorname{Hom}_G(A,E[2]),\\
 \operatorname{Hom}_G(A[1],C[2])=W
       &\longrightarrow\operatorname{Ext}^2_G(C,C)
         \xrightarrow{(-)p}\operatorname{Hom}_G(E,C[2]).
 \end{split}
 \tag{F1R35.9}
\]
Their first images are respectively $\operatorname{im}\nu$
and the shift of $\operatorname{im}\mu$. Thus this projection's
kernel consists exactly of compatible pairs in
$\operatorname{im}\nu\oplus\operatorname{im}\mu$.
On representatives,
\[
                       D(\nu(f),\mu(g))=J(f-g).
 \tag{F1R35.10}
\]
Both representative ambiguities lie in $\ker J$ by
composition, and the restricted $D$ has image exactly
$\operatorname{im}J$. Its kernel consequently has dimension
$r_\nu+r_\mu-r_J$. Combining these two actual projections gives
\[
 \begin{split}
 \operatorname{ext}^2_G(E,E)
 &\ge \operatorname{ext}^2_G(A,A)+\operatorname{ext}^2_G(T,T)
                     -r_D-r_\nu-r_\mu+r_J\\
 &\ge \operatorname{rank}c^2_{z_A}
                  +\operatorname{rank}c^2_{z_T}-3h+r_J.
 \end{split}
 \tag{F1R35.11}
\]
For the last inequality, the original characteristic
evaluation is invariant under the given point lifts:
translations and Picard-zero tensoring fix its $HH^2$
source. Full trace factors the whole-character maps
through these actual invariant endpoint Ext groups.

The $B$-shear of R07 identifies the original source
kernel with $K=K_1\oplus K_2$ by $\bigwedge^2s_B$
and postcomposes full trace by $e^B$. Both maps are
invertible; (F1R35.11) therefore has exactly the same
ranks for the given unmodified objects.
The identity-sector, full-trace and Serre-transpose
maps for the weakly semiregular lift give
\[
                            \operatorname{ext}^2_G(E,E)=18.
 \tag{F1R35.12}
\]
This uses the full balanced character of $E$, not
any endpoint simplicity or semistability assumption.

\paragraph{2. Rationality leaves exactly R26's two rank-two forms.}
Naturality along $i,p$ and (F1R35.9) give the actual
Yoneda memberships
\[
 \operatorname{ev}_A^2(K)\subset\nu(W),\qquad
 \operatorname{ev}_T^2(K)\subset\mu(W).
 \tag{F1R35.13}
\]
Here the displayed kernel is understood through the
actual inverse shear just specified. Their full traces
have rank at most $h\le3$. The rational Hodge-vector
classification in paragraphs 1--3 of R26 is independent
of endpoint ranks. Its explicit polynomial RM projectors,
normalized Mukai pairings and factor-exchanging field
involution therefore apply to these two actual characters.
Their restricted ranks are zero or two and are equal,
since $z_A-z_T=Nv$.

At rank zero, they are $z_A=r_Av$ and $z_T=r_Tv$,
with $r_A,r_T>0$. Each full character map has rank
eighteen. Equations (F1R35.11)--(F1R35.12) would give
$18\ge36-3h\ge27$, a contradiction. Thus the restricted
rank is two and $h\ge2$. This already excludes every
positive-rank coherent adjacent pair with $h\le1$ in
every nonzero virtual rank, without the additional
semistability assumptions of the earlier R19 application.

R26's unique decomposition now gives either
\[
 \begin{split}
 z_T&=a v+b_1v_2+v_1b_2,\\
 z_A&=(a+N)v+b_1v_2+v_1b_2,
 \end{split}
 \qquad 0\ne b_i\perp v_i,
 \tag{F1R35.14}
\]
or the common product in (F1R35.2), initially with
unrestricted $a$. The tensor factors in each alternative
are the exact $F$-defined conjugate Hodge factors of R26.
No cohomological factor is declared to represent a sheaf.

\paragraph{3. The additive alternative is excluded.}
For (F1R35.14), set $\beta_i=\Phi_i^{-1}(b_i)$ and
$P_i=L_{\xi_i}^{-1}L_{\beta_i}$. Each is nonscalar,
since $\beta_i$ is nonzero and independent of $\xi_i$.
The same explicit alternating-form comparison and
Darboux calculation (F1R34.9)--(F1R34.14) apply.
Thus it has two distinct roots with two-dimensional
eigenspaces, or one repeated root with two size-two
Jordan blocks and nonzero square-zero nilpotent of rank two.

Postcompose both actual mixed-character maps by
$(L_{\xi_1}\otimes L_{\xi_2})^{-1}$. They become
\[
              aI+\mathcal S,\qquad (a+N)I+\mathcal S,
       \qquad \mathcal S=P_1\otimes I+I\otimes P_2.
 \tag{F1R35.15}
\]
The two exceptional eigenvalues are $-a$ and $-a-N$,
which are distinct; either may be zero.
For two distinct-root factors, the four tensor summands
have dimension four and their eigenvalues form an additive
rectangle. Its four entries cannot all lie in a two-element
set: if one corner is $u$, both adjacent corners must be
$v$, and the last is $2v-u$, whose equality to $u$ or $v$
forces $u=v$ in characteristic zero. At least four
dimensions therefore contribute to neither kernel.

For one repeated-root factor, each of the other factor's
two eigenspaces gives an eight-dimensional summand.
Its only possible exceptional eigenspace has dimension
four, since the nilpotent is $N_1\otimes I_2$ of rank four.
For two repeated-root factors, each of the four actual
tensor Jordan planes has, after subtracting the sole
eigenvalue, the map
\[
 ee'\longmapsto0,\quad ef'\longmapsto ee',\quad
 fe'\longmapsto ee',\quad ff'\longmapsto ef'+fe'.
 \tag{F1R35.16}
\]
It has rank two and kernel dimension two. Thus in these
last two cases the two exceptional kernels together
have dimension at most eight. In all cases their
combined dimension is at most twelve, and the two
mixed ranks sum to at least twenty.

Each even block has rank two for each endpoint, including
$a=0$ and $a+N=0$: its two coefficient functionals pair
against independent $\xi_i,\beta_i$, and its target
vectors span the independent pair $\xi_j,\beta_j$.
The full character ranks therefore sum to at least
$20+8=28$. Now (F1R35.11), with $r_J\ge0$, gives
$18\ge28-3h\ge19$, a contradiction. This excludes
(F1R35.14) at the entire stated endpoint scope.

\paragraph{4. The product always has a two-dimensional coupling pairing.}
It remains to consider the product in (F1R35.2).
For each block choose an actual $k_i\in K_i$ with
$D_{k_i}p_i=\omega_i\ne0$. Such a choice exists because
$p_i$ is not proportional to $v_i$ and the nondegenerate
pairing on $K_i$ detects its complementary component.
Equation (F1R35.13) supplies $f_i\in W$ with
$\mu(f_i)=\operatorname{ev}_T^2(k_i)$.
The actual BF module identity gives
\[
 \operatorname{Tr}_T\bigl(\mu(f_i)[2]\mu(f_i)\bigr)=0,
 \qquad
 \operatorname{Tr}_T\bigl(\mu(f_1)[2]\mu(f_2)\bigr)
                          =c\omega_1\omega_2\ne0.
 \tag{F1R35.17}
\]
The balanced term is killed in both computations,
and the second application in one block exceeds its
exterior degree. Total degree four leaves precisely
the ordinary Serre trace; both moved degrees are even.
Under invariant CY4 duality, pairing $J(f)$ with $g$
is this same actual trace of $\mu(g)[2]\mu(f)$.
Thus its restriction to the span of $f_1,f_2$ is
nondegenerate and
\[
                                r_J\ge2.
 \tag{F1R35.18}
\]
No value of $a$ was excluded in this argument. In
particular (F1R35.18) holds at all three exceptional
coefficients, and no separate nondegeneracy of a
factor's quadratic norm was assumed.

\paragraph{5. The nonexceptional product coefficients are excluded.}
The actual operators are now
$P_i=L_{\xi_i}^{-1}L_{\alpha_i}$ and
$\mathcal P=P_1\otimes P_2$. They have exactly the
quadratic alternatives proved in R34. The mixed maps
are $aI+c\mathcal P$ and $(a+N)I+c\mathcal P$.

For $a(a+N)(2a+N)\ne0$, their two exceptional values
$-a/c$ and $-(a+N)/c$ are nonzero, distinct and not
opposite. The complete multiplicative-rectangle and
Jordan tensor calculation (F1R34.15)--(F1R34.19)
therefore gives mixed rank sum at least twenty.
This argument requires precisely the three displayed
nonvanishings; no positivity premise is transported
from R28. Both even blocks have rank two for each
endpoint, since $a,a+N,c\ne0$. Thus the full sum is
at least twenty-eight and (F1R35.11), (F1R35.18) give
\[
                    18\ge28-3h+2\ge21,
 \tag{F1R35.19}
\]
a contradiction. Only the three coefficients in
(F1R35.2) remain.

\paragraph{6. The opposite-root coefficient.}
At $a=-N/2$, both even blocks still have rank two
for each endpoint. The two exceptional tensor roots
are $u=N/(2c)$ and $-u$, with $u\ne0$.
If either factor has a repeated root, the explicit
R34 tensor bases give combined exceptional nullity
at most eight. If both have distinct roots but some
product lies outside $\{u,-u\}$, the combined
nullity is at most twelve. Either case again gives
full rank sum at least twenty-eight and is excluded
by (F1R35.19).

Consequently every product of the two pairs of roots
belongs to $\{u,-u\}$. All roots are nonzero. In each
row and column the entries differ, so the two roots
of each factor are negatives of one another. Write
them as $\lambda_i,-\lambda_i$. Their product has
square $u^2$, giving exactly the third equality in
(F1R35.3). The four actual eigenspace tensors split
equally between the exceptional roots. Each mixed
map has rank eight and each whole character map
has rank twelve.

For $h=2$, (F1R35.11) and (F1R35.18) now give
$18\ge24-6+2=20$, still a contradiction. The sole
unexcluded value of $h$ at this coefficient is three,
as stated in the table. No existence assertion is
made for that remaining spectral configuration.

\paragraph{7. A zero balanced coefficient and every zero-root case.}
First take $a=0$. The mixed maps, after multiplication
by nonzero scalars, are $\mathcal P$ and
$\mathcal P-vI$, where $v=-N/c\ne0$.
The pure tensor endpoint has even rank one in each
block, and the other endpoint has even rank two in
each block. Their even rank sum is six.

Each nonscalar $P_i$ has rank two or four. This can
also be seen directly from $L_{\alpha_i}$: a nonzero
two-vector in four dimensions has wedge map of rank
two or four. If both $P_i$ are invertible,
$\ker\mathcal P=0$ and the explicit distinct-root
or Jordan cases give $\dim\ker(\mathcal P-vI)\le8$.

If exactly one $P_i$ is singular, it has either
distinct roots $0,\lambda_i$ with $\lambda_i\ne0$,
or repeated root zero. Its rank is two, so
$\dim\ker\mathcal P=8$. In the repeated-zero case
the tensor operator has only eigenvalue zero and the
other kernel vanishes. In the distinct-root case,
its zero eigenspace contributes no $v$-eigenvectors;
on the other eight-dimensional summand the nonscalar
invertible factor has a $v/\lambda_i$ eigenspace of
dimension at most two, tensored with a two-dimensional
space. Hence the second kernel has dimension at most
four. The combined nullity is at most twelve.

If both $P_i$ are singular, then
$\operatorname{rank}\mathcal P=4$, so its zero
kernel has dimension twelve. A repeated-zero factor
makes the whole tensor nilpotent, leaving no nonzero
exceptional kernel. Otherwise both are semisimple
with roots $0,\lambda_i$ and $\lambda_i\ne0$.
Their only nonzero tensor eigenvalue is
$\lambda_1\lambda_2$, on a four-dimensional summand.
It contributes to the second kernel exactly when
$c\lambda_1\lambda_2=-N$.

Outside this last configuration, the combined nullity
is at most twelve, the mixed ranks sum to at least
twenty and the whole ranks to at least twenty-six.
Equations (F1R35.11), (F1R35.18) would then give
$18\ge26-3h+2\ge19$, a contradiction. The first
row of (F1R35.3) is therefore the complete remaining
configuration at $a=0$. Its actual whole-character
ranks are six and sixteen.

For $a=-N$, the same two maps occur in reversed
order with nonzero exceptional value $N/c$:
they are $\mathcal P$ and $\mathcal P-(N/c)I$.
This is only an interchange of the two linear
character maps; no duality or interchange of the
original sheaves is presumed. The preceding complete
zero-root analysis gives the second row of
(F1R35.3), with equality $c\lambda_1\lambda_2=N$.
Again the whole-character ranks are six and sixteen.

Finally, the rank-two $L_{\alpha_i}$ in these two rows
implies $\alpha_i^2=0$. Explicitly, splitting off a
nonzero alternating pair puts a two-vector into the
form $e_1e_2+d e_3e_4$; its wedge map has rank two
exactly when $d=0$. R26's exact Mukai pairing identity
therefore gives $B_i(p_i,p_i)=0$. Write
$p_i=r_i+\beta_i+s_iA_i$. The positive rank of the
actual pure endpoint is $R=c r_1r_2$, so each $r_i$
is nonzero. The displayed isotropy says
\[
                 2r_i s_i A_i=\beta_i^2,
                  \qquad p_i=r_i\exp(\beta_i/r_i).
 \tag{F1R35.20}
\]
Thus the same whole tensor is $R\exp(\ell)$ with
$\ell=\beta_1/r_1+\beta_2/r_2$. Its degree-two
component divided by $R$ proves directly that
$\ell$ is rational of type $(1,1)$. Undoing the
actual cohomological shear proves (F1R35.4).
It does not identify the coherent endpoint with
a line bundle or a semihomogeneous bundle.

\paragraph{Framework, source scope and remaining actual problem.}
All endpoint maps above are the original equivariant
truncation arrows and the original $\delta$.
The full trace and degree-four pairing use the BF
Hochschild module maps audited in R33, with its exact
graded signs and the explicit R19/R26 exterior-module
isomorphisms. TR3 and the two Hom predecessors are the
Stacks 05QK/0149 maps freshly checked there.
The complex pencil identity and bases are the direct
R34 proof, with its acknowledged classical
Ramakrishna--Ansari--Costa precedent. The additive
tensor calculation and every exceptional product root
were proved explicitly here.

This strictly removes the balanced and additive
whole-character alternatives at all positive endpoint
ranks and removes every nonexceptional product coefficient
for $h\le3$. It also closes $h\le1$ without endpoint
semistability or virtual-rank-one assumptions. The
three rows of (F1R35.3) remain actual existence and
coupling problems; higher $h$ and other amplitudes
retain their earlier unresolved scopes. No row in
that table is represented by a hypothetical object.
